
\documentclass{article}
\usepackage{colortbl}
\usepackage{makecell}
\usepackage{multirow}
\usepackage{supertabular}

\begin{document}

\newcounter{utterance}

\centering \large Interaction Transcript for game `hot\_air\_balloon', experiment `air\_balloon\_survival\_en\_reasoning off\_hard', episode 0 with Qwen3{-}32B{-}t0.0.
\vspace{24pt}

{ \footnotesize  \setcounter{utterance}{1}
\setlength{\tabcolsep}{0pt}
\begin{supertabular}{c@{$\;$}|p{.15\linewidth}@{}p{.15\linewidth}p{.15\linewidth}p{.15\linewidth}p{.15\linewidth}p{.15\linewidth}}
    \# & $\;$A & \multicolumn{4}{c}{Game Master} & $\;\:$B\\
    \hline

    \theutterance \stepcounter{utterance}  
    & & \multicolumn{4}{p{0.6\linewidth}}{
        \cellcolor[rgb]{0.9,0.9,0.9}{
            \makecell[{{p{\linewidth}}}]{
                \texttt{\tiny{[P1$\langle$GM]}}
                \texttt{You are participating in a collaborative negotiation game.} \\
\\ 
\texttt{Together with another participant, you must agree on a single set of items that will be kept. Each of you has your own view of how much each item matters to you (importance). You do not know how the other participant values the items. Additionally, you are given the effort each item demands.} \\
\texttt{You may only agree on a set if the total effort of the selected items does not exceed a shared limit:} \\
\\ 
\texttt{LIMIT = 3683} \\
\\ 
\texttt{Here are the individual item effort values:} \\
\\ 
\texttt{Item effort = \{"A21": 548, "B93": 154, "A11": 60, "B33": 517, "C08": 334, "B39": 542, "B96": 707, "C85": 139, "A82": 661, "B31": 780, "B20": 832, "A60": 913, "B09": 219, "C81": 323, "B21": 638\}} \\
\\ 
\texttt{Here is your personal view on the importance of each item:} \\
\\ 
\texttt{Item importance values = \{"A21": 138, "B93": 583, "A11": 868, "B33": 822, "C08": 783, "B39": 65, "B96": 262, "C85": 121, "A82": 508, "B31": 780, "B20": 461, "A60": 484, "B09": 668, "C81": 389, "B21": 808\}} \\
\\ 
\texttt{Goal:} \\
\\ 
\texttt{Your goal is to negotiate a shared set of items that benefits you as much as possible (i.e., maximizes total importance to YOU), while staying within the LIMIT. You are not required to make a PROPOSAL in every message {-} you can simply negotiate as well. All tactics are allowed!} \\
\\ 
\texttt{Interaction Protocol:} \\
\\ 
\texttt{You may only use the following structured formats in a message:} \\
\\ 
\texttt{PROPOSAL: \{'A', 'B', 'C', …\}} \\
\texttt{Propose keeping exactly those items.} \\
\texttt{REFUSE: \{'A', 'B', 'C', …\}} \\
\texttt{Explicitly reject opponent's proposal.} \\
\texttt{ARGUMENT: \{'...'\}} \\
\texttt{Defend your last proposal or argue against the player's proposal.} \\
\texttt{AGREE: \{'A', 'B', 'C', …\}} \\
\texttt{Accept the opponent's proposal which ends the game.} \\
\\ 
\\ 
\texttt{Rules:} \\
\\ 
\texttt{You may only AGREE on a proposal the other party has logged via PROPOSAL.} \\
\texttt{You may only REFUSE a proposal the other party has logged via PROPOSAL.} \\
\texttt{Total effort of any PROPOSAL or AGREE set must be ≤ LIMIT.} \\
\texttt{Do NOT reveal your hidden importance scores.} \\
\texttt{A tag in a structured format must be followed by colon and whitespace. The argument must be a python set containing 0 or more strings.} \\
\texttt{So, it must be of the form TAG: \{...\}} \\
\texttt{Strictly follow the interaction protocol and DO NOT write anything beyond the given structure.} \\
\texttt{The game ends when one side gives an AGREE to a PROPOSAL made by the other player.} \\
\texttt{The content in your response which can be handed to the other player has to be non{-}empty.} \\
\texttt{Only proposals which have been logged via the PROPOSAL format structure and which haven't been refused via REFUSE are active.} \\
\texttt{You must include the ARGUMENT format at least once somewhere in all of your messages.} \\
            }
        }
    }
    & \\ \\

    \theutterance \stepcounter{utterance}  
    & \multicolumn{4}{p{0.6\linewidth}}{
        \cellcolor[rgb]{0.9,0.9,0.9}{
            \makecell[{{p{\linewidth}}}]{
                \texttt{\tiny{[P1$\rangle$GM]}}
                \texttt{<think>} \\
\texttt{Okay, let's see. I need to negotiate a set of items that maximizes my total importance without exceeding the effort limit of 3683. First, I should calculate the total effort of all items to see if the limit is tight. Let me add up all the effort values.} \\
\\ 
\texttt{Looking at the item efforts: A21 is 548, B93 154, A11 60, B33 517, C08 334, B39 542, B96 707, C85 139, A82 661, B31 780, B20 832, A60 913, B09 219, C81 323, B21 638. Let me sum these up step by step.} \\
\\ 
\texttt{Starting with 548 (A21) + 154 = 702. +60 = 762. +517 = 1279. +334 = 1613. +542 = 2155. +707 = 2862. +139 = 3001. +661 = 3662. +780 = 4442. Wait, that's already over the limit. So the total of all items is way over 3683. So we need to select a subset that's under 3683.} \\
\\ 
\texttt{My goal is to maximize my own importance. Let me list the items with their importance and effort, maybe sort them by importance per unit effort to prioritize high value items.} \\
\\ 
\texttt{Looking at my importance values:} \\
\\ 
\texttt{A21: 138 (effort 548)} \\
\texttt{B93: 583 (154)} \\
\texttt{A11: 868 (60)} \\
\texttt{B33: 822 (517)} \\
\texttt{C08: 783 (334)} \\
\texttt{B39: 65 (542)} \\
\texttt{B96: 262 (707)} \\
\texttt{C85: 121 (139)} \\
\texttt{A82: 508 (6} \\
            }
        }
    }
    & & \\ \\

    \theutterance \stepcounter{utterance}  
    & & \multicolumn{4}{p{0.6\linewidth}}{
        \cellcolor[rgb]{0.9,0.9,0.9}{
            \makecell[{{p{\linewidth}}}]{
                \texttt{\tiny{[P1$\langle$GM]}}
                \texttt{Your response did not contain an argument. You must include the structured format ARGUMENT: \{...\} somewhere in your response. Try again} \\
            }
        }
    }
    & \\ \\

    \theutterance \stepcounter{utterance}  
    & \multicolumn{4}{p{0.6\linewidth}}{
        \cellcolor[rgb]{0.9,0.9,0.9}{
            \makecell[{{p{\linewidth}}}]{
                \texttt{\tiny{[P1$\rangle$GM]}}
                \texttt{<think>} \\
\texttt{Okay, let me try to figure this out. The user wants me to make a proposal that maximizes my total importance without exceeding the effort limit. First, I need to calculate the total effort of all items to see how tight the limit is. Let me add up all the effort values:} \\
\\ 
\texttt{A21: 548, B93: 154, A11: 60, B33: 517, C08: 334, B39: 542, B96: 707, C85: 139, A82: 661, B31: 780, B20: 832, A60: 913, B09: 219, C81: 323, B21: 638.} \\
\\ 
\texttt{Adding them up step by step: 548 + 154 = 702; +60 = 762; +517 = 1279; +334 = 1613; +542 = 2155; +707 = 2862; +139 = 3001; +661 = 3662; +780 = 4442. Wait, that's already over 3683. So the total of all items is way over the limit. Therefore, we need to select a subset.} \\
\\ 
\texttt{My goal is to maximize my own importance. Let me list the items with their importance and effort, then sort them by importance per unit effort. Let's see:} \\
\\ 
\texttt{A11 has high importance (868) with low effort (60). B33: 822 importance, effort 517. B21: 808 importance, effort 638. B31: 780 importance, effort 780. C08: 783 importance, effort 334. B09: 668, effort 219. B93: 583, effort 154. A82: 508, effort 661. C81:} \\
            }
        }
    }
    & & \\ \\

    \theutterance \stepcounter{utterance}  
    & & \multicolumn{4}{p{0.6\linewidth}}{
        \cellcolor[rgb]{0.9,0.9,0.9}{
            \makecell[{{p{\linewidth}}}]{
                \texttt{\tiny{[P1$\langle$GM]}}
                \texttt{Your response did not contain an argument. You must include the structured format ARGUMENT: \{...\} somewhere in your response. Try again} \\
            }
        }
    }
    & \\ \\

    \theutterance \stepcounter{utterance}  
    & \multicolumn{4}{p{0.6\linewidth}}{
        \cellcolor[rgb]{0.9,0.9,0.9}{
            \makecell[{{p{\linewidth}}}]{
                \texttt{\tiny{[P1$\rangle$GM]}}
                \texttt{<think>} \\
\texttt{Okay, let's try this again. The user wants me to make a proposal that maximizes my total importance without exceeding the effort limit. First, I need to calculate the total effort of all items to see how tight the limit is. Let me add up all the effort values:} \\
\\ 
\texttt{A21: 548, B93: 154, A11: 60, B33: 517, C08: 334, B39: 542, B96: 707, C85: 139, A82: 661, B31: 780, B20: 832, A60: 913, B09: 219, C81: 323, B21: 638.} \\
\\ 
\texttt{Adding them up step by step: 548 + 154 = 702; +60 = 762; +517 = 1279; +334 = 1613; +542 = 2155; +707 = 2862; +139 = 3001; +661 = 3662; +780 = 4442. Wait, that's already over 3683. So the total of all items is way over the limit. Therefore, we need to select a subset.} \\
\\ 
\texttt{My goal is to maximize my own importance. Let me list the items with their importance and effort, then sort them by importance per unit effort. Let's see:} \\
\\ 
\texttt{A11 has high importance (868) with low effort (60). B33: 822 importance, effort 517. B21: 808 importance, effort 638. B31: 780 importance, effort 780. C08: 783 importance, effort 334. B09: 668, effort 219. B93: 583, effort 154. A82: 508, effort 661. C81: 3} \\
            }
        }
    }
    & & \\ \\

    \theutterance \stepcounter{utterance}  
    & & & \multicolumn{2}{p{0.3\linewidth}}{
        \cellcolor[rgb]{0.9,0.9,0.9}{
            \makecell[{{p{\linewidth}}}]{
                \texttt{\tiny{[GM$|$GM]}}
                \texttt{abort} \\
            }
        }
    }
    & & \\ \\

    \theutterance \stepcounter{utterance}  
    & & & \multicolumn{2}{p{0.3\linewidth}}{
        \cellcolor[rgb]{0.9,0.9,0.9}{
            \makecell[{{p{\linewidth}}}]{
                \texttt{\tiny{[GM$|$GM]}}
                \texttt{end game} \\
            }
        }
    }
    & & \\ \\

\end{supertabular}
}

\end{document}
